\documentclass[12px]{article}

\title{Lezione 12 Metodi Numerici di Approssimazione}
\date{2026-04-17}
\author{Federico De Sisti}

\input{../../setup.tex}

\begin{document}
	\maketitle
	\newpage
	\subsection{Metodi RK}
	\[
		y(t_{n+1}) = y(t_n) + \int_{t_n}^{t_{n+1}}f(s,y(s))ds
	.\] 
	Introducendo la formula di quadratura
	\[
		y(t_{n+1})\approx y(t_n) + h \sum^{s}_{i=1}b_if(t_n + c_ih,y(t_n + c_i h))
	.\] 
	$b_i, \ \ i = 1,\ldots,s $ sono i pesi\\
	 $c_i, \ \ i = 1,\ldots, s$ sono i nodi\\
	  $s$ è il numero di stadi.
	  Rimane il fatto che dall'altra parte c'è la soluzione valutata in un tempo diverso, riscriviamo nella forma di Volterra.
	  \[
		  y(t_n + c_ih) \approx y(t_n) + h \sum^{s}_{j=1} a_{ij}f(t_n + c_jh,y(t_n + c_j h))
	  .\] 
	  Ma qui compare di nuovo la soluzione calcolata in un tempo diverso, devo "sbrogliare" la ricorsione, noto che due parti sono uguali ma con indice diverso, definisco quindi
	  \[
		  K_i := f(t_n + c_ih, y(t_n) + h \sum_{j=1}^{s}a_{ij}K_j)\ \ \ i = 1,\ldots, s
	  .\] 
	  Ma guardando l'equazione mi accorgo che già da $K_1$ dipendono da se stessi.\\
	  Mi trovo davanti ad un sistema di $s$ equazioni non lineari.\\
	  chiamando $\K = (K_1,\ldots, K_s)\in\R^s$\\
	  $\K = F(\K)$, quindi  $\K$ è un punto fisso di $F$.\\
	  Arriviamo ora allo schema
	   \[
		   K_i = f(t_n + c_i h, u_n + h \sum^{s}_{j=1}a_{ij}K_j)
	  \] 
	  \[
		  u_{n+1} = u_n + h \sum^{s}_{i=1}b_iK_i
	  .\] 
	  Questo è quello che si chiama un $RK$ a $s$ stadi.\\
	  Il metodo è univocamente determinato da $c_i, a_{ij}, b_i$, per identificarli possiamo costruire la tabella di butcher. \\
	  AGGIUNGI FOTO 10 37\\
	  Dobbiamo trovare delle condizioni sui nodi e sui coefficienti epr garantiri che questo metodo sia consistente e zero stabile.\\
	  \subsection{Metodi RK Espliciti}
	  Non vogliamo che $K_i$ dipenda da se stesso o da quelli successivi, questo si traduce con una matrice dei coefficienti  $\{a_{ij}\}$ strettamente triangolare inferiore.\\
	  Posso calcolare  $K_i$ a cascata.\\
	  \subsubsection{Condizione di compatibilità}
	  \[
		  c_i = \sum^{s}_{j=1}a_{ij} \ \ \forall i=1,\ldots,s
	  .\] 
	  i coefficienti $ a_{ij}$ non possono essere arbitrari.\\[5px]
	  \textbf{Motivazione}\\[5px]
	  $
	  \begin{cases}
	  	\dot y (t) = 1\\
		y(0) = y_0
	  \end{cases} \Rightarrow  y(t) = t + cost = t+ y_0$\\
	  Se scrivo  $y(t_{n+1}) = t_{n+1} + y_0 =t_n + h + y_0= y(t_n) + h$\\
	  \[u_{n+1} = u_n + h \sum^{s}_{i=1}b_iK_i\], con \[K_i = f(t_n+c_ih,u_n + \sum^{s}_{j=1}a_{ij}K_j)=1\]
	  \[
		  u_{n+1} = u_n + h \sum^{s}_{i=1}b_i
	  .\] 
	  Confrontando con la soluzione scopriamo che $ \sum^{s}_{i=1}b_i = 1$, chiaramente sono dei pesi di quadratura, la somma deve pesare $1$.
	  \[
		  \begin{cases}
	  \dot y(t) = f(t,y(t))\\
	  y(0) = y_0
		  \end{cases} \Rightarrow \begin{cases}
		  	\dot t(s) = 1\\
			\dot y(s) = f(s,y(s))\\
			y(0)= 0\\
			t(0) = 0
		  \end{cases} \Rightarrow \begin{cases}
		  	Y(s) =(t(s),y(s))\\
			F(Y(s)) = (1,f(s,y(s)))
		  \end{cases}
	  .\] 
	  Arriviamo al sistema autonomo
	  \[
		  \begin{cases}
	  \dot Y(s) = F(Y(s))\\
	  Y(0) = Y_0 = (0,y_0)
		  \end{cases}
	  .\] 
	  Proviamo a riscrivere il sistema di prima in forma autonoma
	  \[
	  \begin{cases}
	  	\dot y = 1\\
		y(0) = y_0
	  \end{cases} \Rightarrow \begin{cases}
	  	\dot t(s)= 1\\
		\dot y(s) = f(s,y(s)
	  \end{cases}
	  .\] 
	  abbiamo
	  \[
		  U_{n+1} =U_n + h \sum^{s}_{i=1}b_iK_i
	  .\] 
	  \[
		  K_i = F(t_n,c_ih,U_n + \sum^{s}_{j=1}a_{ij}K_j)  = \matrice {1\\f(t_n+c_ih, \sum_j a_{ij}k_j)}
	  .\] 
	  quindi i $K_i$ sono dei vettori.\\
	  Usando ora l'approssimazione
	  \[
		  y(t_n + c_ih) \approx y(t_n) + h \sum^{s}_{j=1}a_{ij}f(t_n+c_jh,y(t_n,c_jh))
	  .\] 
	  \[
	  \begin{cases}
	  	\dot y(t) = 1\\
		y(0) = y_0
	\end{cases} \Rightarrow y(t_{n+1})= y_0 + t_{n+1} = y(t_n) + h
	  .\] \\
	  $y(t_n+c_ih) = t_n + c_ih + y_0$\\
	  Valutando la prima approssimazione
	  \[
		  y(t_n) + h \sum^{s}_{j=1}a_{ij} = t_n + y_0 + \sum^{s}_{j=1}a_{ij} \Rightarrow c_i = \sum^{s}_{j=1}a_{ij}
	  .\] 
	  confrontandola con l'equazione sopra
	  \begin{prop}
	  Un metodo $RK$ è consistente $ \Leftrightarrow \sum^{s}_{i=1}b_i = 1$
	  \end{prop}
	  \begin{dimo}
	  \[
		  y_{n+1} - \left(y_n + h \sum^{s}_{i}b_iK_i \right) = h \tau_{n+1}(h)
	  .\] 
	  \[
		  K_i = f(t_n + c_i h , y_n + h \sum^{s}_{j=1}a_{ij}K_j \xrightarrow{h \rightarrow 0} f(t_n,y_n) + O(h)
	  .\] 
	  \[
		  y_{n+1} - (y_n + h \sum^{s}_{i=1}b_i(f(t_n,y_n) + O(h)) = h\tau_{n+1}(h)
	  .\] 
	  \[
		  y_{n+1} = y_n + h\dot y_n + O(h^2) = y_n + hf(t_n,y_n) + O(h^2)
	  .\] 
	  \[
		  \cancel{y_n} + hf(t_n,y_n) - \cancel{y_n} - h \sum^{s}_{i=1}f(t_n,y_n) - O(h^2) \sum^{s}_{i=1}b_i = h\tau_{n+1}(h)
	  .\] 
	  \[
		  hf(t_n,y_n) (1- \sum^{s}_{i=1}b_i) - O(h^2) \sum^{s}_{i=1}b_i = j\tau_{n+1}(h)
	  .\] 
	  \[
		  \Rightarrow \tau_{n+1}(h) = f(t_n,y_n)(1 - \sum^{s}_{i=1}b_i) + O(h) \xrightarrow{h \rightarrow 0}0 \Leftrightarrow \sum^{s}_{i}b_i = 1
	  .\] 
  \end{dimo}
	  \textbf{Esempio}\\
	  1 Stadio esplicito $\matrice{c_1 & a_{11}\\ & b_1}$ \ \ $\matrice{0 & 0 \\ & 1}$ \\
	  questo è Eulero in avanti  $k_1 = f(t_n+c_1h,u_n+h \sum^{s}_{j=1}a_{ij}k_j) \Rightarrow k_1 = f(t_n,u_n)$ \\
	  $u_{n+1} = u_n + h \sum^{s}_{i}b_ik_i = u_n + h b_1 k_1 = u_n + h f(t_n,u_n)$\\
	  \textbf{Idea per costruire un metodo RK a s stadi}\\
	  Vogliamo massimizzare l'ordine di consistenza al variare del numero di stadi.\\
	  Bisogna fare lo sviluppo di Taylor e si cerca di annullare più derivate possibili.\\
	  2 Stadi:\\
	  $\matrice{0 & 0 & 0\\
		  c_2 & a_{21}& 0\\
		      & b_1 & b_2} \Rightarrow \matrice{0 & 0 & 0\\ c_2 & c_2 &0\\
		      & b_1 &b_2} $ e si deve avere $b_1 +b_2 = 1$\\
 siccome la prima fila è nulla\\
		      $u_{n+1} = u_n + hb_1K_1 + hb_2K_2 = u_n + hb_1f(t_n,u_n)+h b_2f(t_n+c_2h, u_n + hc_2f(t_n,u_n))$\\
		      $y_{n+1} = y_n + hb_1 f(t_n,y_n) + h b_2 (\circledast)+h\tau_{n+1}(h)$\\
Dobbiamo fare uno sviluppo di taylor
\[
	\circledast f(t_n + c_2h , y_n + hc_2f(t_n,u_n)) 
\] 
\[
=f(t_n,y_n) +[ \frac {\partial}{\partial t} f(t_n,y_n)c_2 + \frac{\partial f}{\partial y}(t_n,y_n)c_2 f(t_n,y_n)]h + O(h^2)
.\] 
\[
	y_{n+1} = y_n + hb_1f_n + hb_2(f_n + c_2h(\frac{\partial}{\partial t} + \frac{\partial f_n}{\partial y}f_n) +O(h^2)) + h\tau_{n+1}(h)
.\] 
Taylor su $y_{n+1}$ calcolato in $t_n$
 \[
	 y_{n+1} = y_n + h\dot y_n+ \frac{h^2}{2}\ddot y_n + O(h^3)
.\] 
\[
\dot y(t) = f(t,y(t))
.\] 
\[
	\ddot y(t) = \frac{\partial f}{\partial t}(t,y(t)) + \frac{\partial f}{\partial y}(t,y(t))\dot y(t) = \frac{\partial f}{\partial t} + \frac{\partial f}{\partial y}f
.\] 
Possiamo quindi riscrivere
 \[
	 y_{n+1} = y_n + hf_n + \frac{h^2}2 (\frac{\partial f_n}{\partial t} + \frac{\partial f_n}{\partial y} f_n) + O(h^3)
.\] 
Possiamo quindi scrivere e semplificare
\[
	\frac{h^2}2 \left(\frac{\partial f_n}{\partial t} + \frac{\partial f_n}{\partial y}f_n \right) + O(h^3) + h b_2c_2 \left(\frac{\partial}{\partial t}f_n + \frac{\partial f_n}{\partial y}f_n \right) + hb_2o(h^2) + h\tau{n+1}(h)
.\] 
i due coefficienti che ci interessano si cancellano se $b_2c_2=\frac 12 \Leftrightarrow h\tau_{n+1}(h) = O(h^3) \Rightarrow  \tau_{n+1}(h) = O(h^2)$ \\
Tutti i metodi $RK2$ esplici sono della forma
 \[
\begin{cases}
	b_1 + b_2 = 1\\
	c_2 b_2 = \frac 12
\end{cases} \Rightarrow  \begin{cases}
	b_1 = 1 - \alpha, \ \ b_2 = \alpha\\
	c_2 = \frac{1}{2b_2} = \frac 1 {2\alpha}
\end{cases}
.\] 
quindi tutti i metodi sono della famiglia con questa tabella $\matrice{\frac 1 {2\alpha} & \frac 1 {2\alpha} & 0\\
											 & 1 - \alpha & \alpha}$\\
 \textbf{Esempio}\\
 $\alpha = \frac 12 , \  \b_1 = \frac 12, \ \ b_2 = \frac 12, \ \ c_2 = 1$\\
 è il metodo di Heun
  \[
	  u_{n+1} = u_n + \frac h 2 (f_n + f(t_n+ h, u_n+hf_n))
 .\] 
 $\alpha = 1, \ \ b_1 = 0, \ \ b_2 = 1, \ \ c_2 = \frac 12$
 \[
	 u_{n+1} = u_n + h f (t_n + \frac h 2, u_n + \frac h 2 f_n)
 .\] 
 Eulero modificato, che è di ordine 2! Fa un mezzo passo di Eulero e calcola la dinamica in mezzo.

	
\end{document}
